% MOMENTBALANCE3
\documentclass[11pt]{article}
\usepackage[margin=.88in]{geometry}
\usepackage{amsmath,amssymb,amsthm}
\newtheorem{proposition}{Proposition}
\begin{document}

\begin{center}
{\Large MOMENTBALANCE3: Exact pair-moment balancing and held-out cubic tests}\\[4pt]
\textbf{Date: 2026-10-04}\\
\textbf{Status: COMPUTATIONAL + elementary moment lemma}\\
\textbf{P versus NP status: OPEN}
\end{center}

\section{Objective}

ORDER3EXACT1 suggested that, after pair matching, no reproducible exact
third-order orbit distinguisher had yet appeared.  Nearest-neighbor matching,
however, only approximately aligns individual pair vectors.

MOMENTBALANCE3 instead modifies the negative SAT distribution by exponential
tilting so that, on training data, all eight independent signed pair-orbit
means agree exactly with those of the rare UNSAT distribution.

The eight balanced coordinates are
\[
(1,0),(1,1),(2,0),(2,1),(2,2),(3,1),(3,2),(3,3).
\]
The omitted \((0,0)\) coordinate is determined by the fixed total number
\(\binom t2\) of clause pairs.

\section{Method}

Let \(Z(F)\in\mathbb R^8\) be the SAT-standardized pair-orbit vector.
For training SAT formulas \(F\), define
\[
w_\lambda(F)
=
\frac{\exp(\lambda\cdot Z(F))}
{\sum_{G\in SAT_{train}}\exp(\lambda\cdot Z(G))}.
\]
The vector \(\lambda\) is chosen to solve
\[
\sum_{F\in SAT_{train}}w_\lambda(F)Z(F)
=
\frac1{|UNSAT_{train}|}
\sum_{F\in UNSAT_{train}}Z(F).
\tag{MB}
\]

Every reported training solve reached maximum coordinate error at most
\[
4.4\times10^{-11}.
\]

The learned \(\lambda\) is then frozen and applied to held-out SAT formulas.
No holdout UNSAT formula is used to fit the weights.

\section{Runtime decomposition rule}

Following the project-wide runtime rule introduced this run, every Python
invocation was designed to finish in roughly 30 seconds or less.  Different
values of \(q\) and replicate seeds were run as independent chunks.
The largest sampling chunk, \(q=24,N=16000\), completed in 14.1 seconds.
The reusable script refuses \(N>16000\) and aborts near 25 seconds so a
larger job must be decomposed rather than given a longer timeout.

\section{Primary results}

At density \(t=3.5q\):

\[
\begin{array}{c|c|c|c|c|c}
q&N&U_{train}&U_{hold}&ESS_{hold}&
\max |\Delta_{pair}|_{hold}\;({\rm SD})\\
\hline
16&10000&154&151&1249.6&0.197\\
18&12000&188&178&742.1&0.177\\
20&14000&171&173&1483.2&0.158\\
24&16000&143&130&1842.6&0.224
\end{array}
\]

The five predeclared cubic statistics were:
\[
\sum_i\left[\binom{d_i^+}{2}d_i^-+d_i^+\binom{d_i^-}{2}\right],
\]
\[
\sum_i\left[\binom{d_i^+}{3}+\binom{d_i^-}{3}\right],
\qquad
\sum_i\binom{d_i}{3},
\qquad
\sum_i d_i^2,
\qquad
\sum_i(d_i^+d_i^-)^2.
\]

Their maximum held-out standardized mean differences were
\[
0.069, 0.054, 0.192, 0.133
\]
for \(q=16,18,20,24\), respectively.

The sign of the residual was not stable.  In particular, three independent
\(q=20\) chunks produced positive, positive, and negative cubic residual
vectors.  A second \(q=24\) chunk produced a vector close to zero rather
than reproducing the first negative vector.

\section{Pooled reconciliation}

Pooling the independent chunks while preserving train/holdout separation gives

\[
q=20:
\qquad
U_{train}=416,quad U_{hold}=398,
\]
\[
ESS_{hold}=3499.0,
\qquad
\max|\Delta_{pair}|=0.126\;{\rm SD},
\]
and the five cubic residuals
\[
(0.056,0.047,0.063,0.039,0.073).
\]

For \(q=24\):
\[
U_{train}=249,quad U_{hold}=217,
\]
\[
ESS_{hold}=3039.5,
\qquad
\max|\Delta_{pair}|=0.126\;{\rm SD},
\]
with cubic residuals
\[
(-0.056,-0.082,-0.077,-0.090,-0.040).
\]

Thus after explicit pair-moment balancing, the tested cubic local statistics
are small compared with the original raw pair leakage and do not show a
replicated direction across independent chunks.

\section{Moment interpretation}

If two exact-layer distributions match every squarefree moment
\[
\mathbb E\prod_{i\in S}z_i
\]
for \(|S|\le d\), then they have equal expectation for every multilinear
polynomial of degree at most \(d\).  This elementary fact is the exact
algebraic meaning of the moment-camouflage hierarchy.

MOMENTBALANCE3 does not yet match every individual degree-two monomial.
It matches the complete orbit-averaged degree-two moment vector under the
signed-permutation symmetry class.  Because the target distributions are
intended to be symmetry-invariant, this is the correct invariant quotient,
but the empirical reweighted sample is only approximately symmetric at
finite size.

\section{What this does not prove}

The experiment does not prove:
\begin{itemize}
\item exact degree-three moment equality;
\item asymptotic pair or cubic camouflage;
\item indistinguishability by general monotone circuits;
\item any minimum-MUS lower bound;
\item \(P\ne NP\).
\end{itemize}

The missing load-bearing theorem remains a circuit-to-local-structure
approximation theorem strong enough for general monotone DAGs on this slice.

\section{Next target}

The next experiment is \textbf{MOMENTBALANCE3-ORBIT}: use the same frozen
pair-moment tilt, then compute the complete canonical degree-three orbit vector
on a bounded number of independently chunked formulas at \(q=16,18,20\).
The job must be decomposed so every individual Python invocation remains below
roughly 30 seconds.

If the exact order-three vector remains quiet under moment balancing, then the
next mathematical target is a SAT-specific approximation/sunflower lemma;
otherwise the stable degree-three orbit becomes the next hostile distinguisher.

\section{Status}

Pair-moment balancing: \textbf{COMPUTATIONAL}.  Moment-camouflage linearity
lemma: \textbf{PROVED/EXACT}.  Asymptotic degree-three camouflage and the
general monotone lower bound remain \textbf{OPEN}.  Therefore \(P\ne NP\)
remains OPEN.

\end{document}
